
\subsection{滤过 $A_{n,K}$ 结构的构造}
\label{subsec:7-2-7}

现在用命题~\ref{prop:7-2-35}
在 $L$ 上构造滤过 $A_{n,K}$ 代数。
具体地，要证明命题 7.2.78。
先置
\begin{equation}
\mathfrak m^{\mathrm{geo}}_{k,\beta}
(P_1,\ldots,P_k)
=\bigl(
\mathcal M^{\mathrm{main},\mathfrak d}_{k+1}
(L;\beta;P_1,\ldots,P_k)
^{\mathfrak s_{\mathfrak d,k,\beta,\vec P}},
\operatorname{ev}_0\bigr),
\tag{7.2.77}\label{eq:7-2-77}
\end{equation}
右端含义如 (7.2.36.3)。

\begin{proposition}{7.2.78}\label{prop:7-2-78}
\liftIndex{I254}\liftIndex{I255}%
对任意 $g_0,n,K$，
存在 $g_1>g_0$，
以及 $C_{(g_1)}(L;\Lambda_{0,\mathrm{nov}})$ 上的
滤过 $A_{n,K}$ 结构 $\mathfrak m_{k,\beta}$，
使每当 $P_i\in\mathcal X_{(g_0)}(L)$ 时，
都有
\[
\mathfrak m_{k,\beta}(P_1,\ldots,P_k)
=\mathfrak m^{\mathrm{geo}}_{k,\beta}
(P_1,\ldots,P_k).
\]
\end{proposition}

命题~\ref{prop:7-2-78}
由命题~\ref{prop:7-2-35}
通过纯代数论证推出。
后者所用的数 $\mathcal K$ 将选得只依赖于
$g_0,n,K$。

\begin{remark}{7.2.79}\label{rem:7-2-79}
只要 $A_{n,K}$ 关系两端都有定义，
$\mathfrak m^{\mathrm{geo}}_{k,\beta}$ 就满足该关系。
然而，若 $P_i\in\mathcal X_{g_i}(L)$ 且
\[
\max\{g_i\mid i=1,\ldots,k\}
+\|\beta\|+k-1>\mathcal K,
\]
则 $\mathfrak m^{\mathrm{geo}}_{k,\beta}
(P_1,\ldots,P_k)$ 可能无定义。
这就是需要把
$\mathfrak m^{\mathrm{geo}}_{k,\beta}$
改造成 $\mathfrak m_{k,\beta}$ 的原因。
注~\ref{rem:7-2-88} 还将给出另一种方法。
\end{remark}

\begin{proof}[命题~\ref{prop:7-2-78} 的证明]
沿用第 5.4 节构造典范模型的思路。
先取常数 $c_1$，使对每个
$(\beta,k)\lesssim(n,K)$，均有
\begin{equation}
\|\beta\|+k<c_1,
\tag{7.2.80}\label{eq:7-2-80}
\end{equation}
并置 $g_m=g_0+mc_1$。
稍后将在证明中固定另一个正整数 $M$，
使 $g_M\le\mathcal K$。
注意 $C_{(g_m)}(L;\mathbb Q)$ 是由
$\mathcal X_{(g_m)}(L)$ 生成的 $\mathbb Q$-向量空间。
考虑链映射
\[
\Pi_{C_{(g_1)}(L;\mathbb Q)}\colon
C_{(g_M)}(L;\mathbb Q)
\longrightarrow C_{(g_M)}(L;\mathbb Q),
\]
以下简记为 $\Pi$；要求
\[
\Pi^2=\Pi,
\qquad
\operatorname{Im}\Pi=C_{(g_1)}(L;\mathbb Q).
\]

\begin{lemma}{7.2.81}\label{lem:7-2-81}
存在这样的 $\Pi$，以及一列 $\mathbb Q$-线性映射
\[
\mathfrak G_k\colon
C^k_{(g_M)}(L;\mathbb Q)
\longrightarrow C^{k-1}_{(g_M)}(L;\mathbb Q),
\]
满足以下性质：

\noindent\textup{(7.2.82.1)}\quad
在 $C^k_{(g_M)}(L;\mathbb Q)$ 上，
$1-\Pi=-\overline{\mathfrak m}_1\circ\mathfrak G_k
-\mathfrak G_{k+1}\circ\overline{\mathfrak m}_1$。

\noindent\textup{(7.2.82.2)}\quad
$\mathfrak G_k\circ\mathfrak G_{k+1}=0$。

\noindent\textup{(7.2.82.3)}\quad
$\mathfrak G_k$ 保持
$C_{(g_M)}(L;\mathbb Q)$ 上的滤过
$\{C_{(g_m)}(L;\mathbb Q)\}_m$；
即
\[
\mathfrak G_k(C_{(g_m)}(L;\mathbb Q))
\subseteq C_{(g_m)}(L;\mathbb Q).
\]

\noindent\textup{(7.2.82.4)}\quad
$\mathfrak G_k$ 在
$C_{(g_1)}(L;\mathbb Q)$ 上为零。
\end{lemma}

\begin{proof}[引理~\ref{lem:7-2-81} 的证明]
可取 $A_m\subset C_{(g_m)}(L;\mathbb Q)$，使
\[
C_{(g_m)}(L;\mathbb Q)
=C_{(g_{m-1})}(L;\mathbb Q)
\oplus A_m\oplus\overline{\mathfrak m}_1(A_m).
\]
于是得到如下逐级直和分解；原书下一式的求和界有一处误植。%
\translatornote{原书本式把直和下、上界误写为
$m=g_1+1$ 与 $g_M$；
$A_m$ 的下标是 $g_m$ 的序号，
由紧前的逐级分解应从 $m=2$ 求至 $M$。
参见疑误 A-0439。}
\[
C_{(g_M)}(L;\mathbb Q)
=C_{(g_1)}(L;\mathbb Q)
\oplus\bigoplus_{m=2}^{M}
\bigl(A_m\oplus\overline{\mathfrak m}_1(A_m)\bigr).
\]
据此在各直和项上定义
\[
\Pi(x)=
\begin{cases}
0,&x\in A_m,\\
x,&x\in C_{(g_1)}(L;\mathbb Q),\\
0,&x\in\overline{\mathfrak m}_1(A_m),
\end{cases}
\]
以及
\[
\mathfrak G(x)=
\begin{cases}
0,&x\in A_m,\\
0,&x\in C_{(g_1)}(L;\mathbb Q),\\
-(\overline{\mathfrak m}_1|_{A_m})^{-1}(x),
  &x\in\overline{\mathfrak m}_1(A_m).
\end{cases}
\]
容易看出，$\Pi$ 和 $\mathfrak G$
具有要求的性质。
\end{proof}

\begin{lemma}{7.2.83}\label{lem:7-2-83}
若 $x_i\in C_{(g_m)}(L;\mathbb Q)$，
$i=1,\ldots,k$，
$(\beta,k)\lesssim(n,K)$，且 $g_m\le\mathcal K$，
则
\[
\mathfrak m^{\mathrm{geo}}_{k,\beta}
(x_1,\ldots,x_k)
\in C_{(g_{m+1})}(L;\mathbb Q).
\]
\end{lemma}
这是命题~\ref{prop:7-2-35}
以及 $c_1,g_i$ 的选法的直接推论。

现在仿照第 5.4 节构造典范模型，
以 $\Pi$ 代替 $\Pi_H$，
并以引理~\ref{lem:7-2-81}
代替引理 5.4.28。
具体地，对 $(\beta,k)\lesssim(n,K)$
定义 $\mathfrak i_{k,\beta}$、
$\mathfrak m_{k,\beta}$。
沿用第 5.4.4 小节的记号：
设
$\Gamma=(T,i,v_0,V_{\mathrm{tad}},\eta)
\in\mathcal G^+_{k+1}$，并置
\[
\beta(\Gamma)=
\sum_{v\in C^0_{\mathrm{int}}(T)}\beta_{\eta(v)}.
\]
将 $\mathfrak G(L)$ 的元素按能量非递减次序
列为 $\{\beta_0,\beta_1,\ldots\}$，
即 $E(\beta_i)\le E(\beta_{i+1})$。
对满足
$(\beta(\Gamma),k)\lesssim(n,K)$ 的
$\Gamma\in\mathcal G^+_{k+1}$，
归纳定义 $\mathfrak i_\Gamma$ 与
$\mathfrak m_\Gamma$：
\begin{equation}
\left\{\begin{aligned}
\mathfrak m_\Gamma(\mathbf x)
&=\sum_a\Pi\circ
\mathfrak m^{\mathrm{geo}}_{\ell,\beta_{\eta(v_1)}}
\bigl(\mathfrak i_{\Gamma_1}(\mathbf x_a^{(1)})
\otimes\cdots\otimes
\mathfrak i_{\Gamma_\ell}(\mathbf x_a^{(\ell)})\bigr),\\
\mathfrak i_\Gamma(\mathbf x)
&=\sum_a\mathfrak G\circ
\mathfrak m^{\mathrm{geo}}_{\ell,\beta_{\eta(v_1)}}
\bigl(\mathfrak i_{\Gamma_1}(\mathbf x_a^{(1)})
\otimes\cdots\otimes
\mathfrak i_{\Gamma_\ell}(\mathbf x_a^{(\ell)})\bigr).
\end{aligned}\right.
\tag{7.2.84}\label{eq:7-2-84}
\end{equation}
这里的 $\Gamma_i$ 如式 (5.4.36)；
那里用记号 $\mathfrak f$，
此处改用 $\mathfrak i$，
因为第 7.2.9 小节还需以 $\mathfrak f$
表示另一个 $A_{n,K}$ 同态。
另置
\begin{equation}
\mathfrak m_{\Gamma_{k+1,\eta}}
=\Pi\circ\mathfrak m^{\mathrm{geo}}_{k,\beta_{\eta(v)}},
\qquad
\mathfrak i_{\Gamma_{k+1,\eta}}
=\mathfrak G\circ
\mathfrak m^{\mathrm{geo}}_{k,\beta_{\eta(v)}},
\tag{7.2.85}\label{eq:7-2-85}
\end{equation}
以启动归纳；这里的记号如式 (5.4.34)。

须注意 $\mathfrak m^{\mathrm{geo}}$ 只部分定义。
因此要证明，当
$\mathbf x\in B_kC_{(g_1)}(L;\mathbb Q)$ 时，
式~\eqref{eq:7-2-84} 的右端良定义。
为此证明引理 7.2.86。
对
$\Gamma=(T,i,v_0,V_{\mathrm{tad}},\eta)
\in\mathcal G^+_{k+1}$，
置
\[
m(\Gamma)=\#\{\Gamma'\mid\Gamma'<\Gamma\},
\]
其中 $<$ 如式 (5.4.33)。

\begin{lemma}{7.2.86}\label{lem:7-2-86}
若
$\mathbf x\in B_kC_{(g_m)}(L;\mathbb Q)$、
$g_{m+m(\Gamma)}<\mathcal K$，且
$\ell+\eta(v_1)<\mathcal K$，
则
\[
\mathfrak m^{\mathrm{geo}}_{\ell,\beta_{\eta(v_1)}}
\bigl(\mathfrak i_{\Gamma_1}(\mathbf x_a^{(1)})
\otimes\cdots\otimes
\mathfrak i_{\Gamma_\ell}(\mathbf x_a^{(\ell)})\bigr)
\]
良定义，且属于
$C_{(g_{m+m(\Gamma)})}(L;\mathbb Q)$。
此外，若 $k+\|\beta(\Gamma)\|\le\mathcal K$
且 $\Gamma\in\mathcal G^+_{k+1}$，
则 $\mathfrak m_\Gamma(\mathbf x)$、
$\mathfrak i_\Gamma(\mathbf x)$ 良定义，
并有
\[
\mathfrak i_\Gamma(\mathbf x)
\in C_{(g_{m+m(\Gamma)})}(L;\mathbb Q).
\]
\end{lemma}

\begin{proof}
第二部分由第一部分、式~\eqref{eq:7-2-84}
及 (7.2.82.3) 得到。
从 $\Gamma=\Gamma_{k+1,\eta}$ 开始归纳；
其定义见式 (5.4.33) 后的段落。
此时第一部分由引理~\ref{lem:7-2-83} 得到。

假设结论已对所有 $\Gamma'<\Gamma$ 成立，
再证 $\Gamma$ 的情形。
令 $\Gamma_i$ 如式~\eqref{eq:7-2-84}。
引理 5.4.35 给出 $\Gamma_i<\Gamma$；
因此由归纳假设，
\[
\mathfrak i_{\Gamma_i}
(\mathbf x_a^{(i)})
\in C_{(g_{m+m(\Gamma_i)})}(L;\mathbb Q)
\subset C_{(g_{m+m(\Gamma)-1})}(L;\mathbb Q).
\]
再由引理~\ref{lem:7-2-83}，
\[
\mathfrak m^{\mathrm{geo}}_{\ell,\beta_{\eta(v_1)}}
\bigl(\mathfrak i_{\Gamma_1}(\mathbf x_a^{(1)})
\otimes\cdots\otimes
\mathfrak i_{\Gamma_\ell}(\mathbf x_a^{(\ell)})\bigr)
\in C_{(g_{m+m(\Gamma)})}(L;\mathbb Q).
\]
\end{proof}

现在可固定整数 $M$。取
\[
M=\sup\left\{m(\Gamma)\ \middle|\
\begin{gathered}
k=0,1,2,\ldots,\quad
\Gamma\in\mathcal G^+_{k+1},\\[-2pt]
(\beta(\Gamma),k)\lesssim(n,K)
\end{gathered}\right\}+1,
\]
\translatornote{源书此处未说明 $\mathcal K$ 的选取；为使下文
$g_M\le\mathcal K$ 的叙述自明，译文补一句，见疑误 A-0753。}
并把 $\mathcal K$ 选得足够大，使
$g_M\le\mathcal K$。
由引理~\ref{lem:7-2-86}，
若
$\mathbf x\in B_kC_{(g_1)}(L;\mathbb Q)$、
$\Gamma\in\mathcal G^+_{k+1}$ 且
$(\beta(\Gamma),k)\lesssim(n,K)$，
则 $\mathfrak m_\Gamma(\mathbf x)$、
$\mathfrak i_\Gamma(\mathbf x)$
均良定义。因此置
\begin{equation}
\mathfrak m_{k,\beta}(\mathbf x)
=\sum_{\substack{\Gamma\in\mathcal G^+_{k+1}\\
                  \beta(\Gamma)=\beta}}
\mathfrak m_\Gamma(\mathbf x).
\tag{7.2.87}\label{eq:7-2-87}
\end{equation}
由引理~\ref{lem:7-2-86}，
右端良定义；按定义有
$\mathfrak m_{k,\beta}(\mathbf x)
\in C_{(g_1)}(L;\mathbb Q)$。
与定理 5.4.2 的证明相同，
可证 $\mathfrak m_{k,\beta}$ 满足
$A_{n,K}$ 关系。

若 $\mathbf x\in B_kC_{(g_0)}(L;\mathbb Q)$，
则
$\mathfrak m^{\mathrm{geo}}_{k,\beta_{\eta(v)}}
(\mathbf x)\in C_{(g_1)}(L;\mathbb Q)$。
由 (7.2.82.4) 得
$\mathfrak i_{\Gamma_{k+1,\eta}}(\mathbf x)=0$。
以此为起点归纳，可证
$\mathfrak i_\Gamma(\mathbf x)=0$；
于是按式~\eqref{eq:7-2-87}，
对每个 $\beta$ 都有下式。%
\translatornote{源书末句写成无 $\beta$ 下标的
$\mathfrak m_k(\mathbf x)
=\mathfrak m_k^{\mathrm{geo}}(\mathbf x)$；
命题陈述及上述树形归纳均逐个 $\beta$ 处理，
此处按分量式呈现。参见疑误 A-0440。}
\[
\mathfrak m_{k,\beta}(\mathbf x)
=\mathfrak m^{\mathrm{geo}}_{k,\beta}(\mathbf x).
\]
命题~\ref{prop:7-2-78} 证毕。
\end{proof}

证明过程中还得到下列结果。
\begin{proposition}{7.2.78bis}\label{prop:7-2-78bis}
在命题~\ref{prop:7-2-78} 的情形下，
存在同态 $\mathfrak i_{k,\beta}$，
使对
$(\beta,k)\lesssim(n,K)$、
$\mathbf x\in B_kC_{(g_1)}(L;\mathbb Q)$
有恒等式
\begin{equation*}
\begin{aligned}
&\sum_{\substack{m,i,\ \beta'+\beta_1+\cdots+\beta_m=\beta\\
                   k_1+\cdots+k_m=k}}
\mathfrak m^{\mathrm{geo}}_{m,\beta'}
\bigl(\mathfrak i_{k_1,\beta_1}(\mathbf x_i^{(1)}),
\ldots,
\mathfrak i_{k_m,\beta_m}(\mathbf x_i^{(m)})\bigr)\\
&\quad=
\sum_{\substack{\beta_1+\beta_2=\beta\\
                   k_1+k_2=k+1}}
\sum_i(-1)^{\deg'\mathbf x_i^{(1)}}
\mathfrak i_{k_2,\beta_2}
\bigl(\mathbf x_i^{(1)},
\mathfrak m_{k_1,\beta_1}(\mathbf x_i^{(2)}),
\mathbf x_i^{(3)}\bigr).
\end{aligned}
\end{equation*}
此外，若 $(\beta,k)\ne(\beta_0,1)$，
则 $\mathfrak i_{k,\beta}$ 在
$B_kC_{(g_0)}(L;\mathbb Q)$ 上为零。
\end{proposition}
这里 $\mathfrak i_{k,\beta}$ 的定义域
见引理~\ref{lem:7-2-86}；
上式两端的良定义性也是命题的一部分。

\begin{remark}{7.2.88}\label{rem:7-2-88}
本小节通过代数修正命题~\ref{prop:7-2-35}
给出的运算，得到滤过 $A_{n,K}$ 代数。
另一条途径是如下修改该命题的陈述与证明。

命题~\ref{prop:7-2-35}
的证明只进行有限步归纳，
利用 (7.2.56) 中归纳步数的有限性，
分别控制扰动模空间与多值截面
距原对象的距离。
为此，当时只构造
$\mathcal M_{k+1}(L;\beta;\vec P)$
在 $k+\|\beta\|$ 及各 $P_i$ 的代
均受 $\mathcal K$ 限制时的扰动。

也可以修改 (7.2.56) 的条件，
使多值截面的构造可针对无限多个 $P_i$ 重复，
只要求 $(k,\beta)$ 满足
$k+\|\beta\|\le\mathcal K$。
为此把 (7.2.56.1) 改为下面的条件；
其中第二句比较的是原映射的零集。%
\translatornote{源书第二句写“原仓西映射的邻域”，
但扰动模空间是零集；参照命题 7.2.35
的 (7.2.36.5) 与 (7.2.56.1)，
这里应与原映射的零集比较。
参见疑误 A-0441。}
此条件可通过归纳保证：

\noindent\textup{(*)}\quad
$\mathcal M_{k+1}(L;\beta;\vec P)$
上的多值截面
$\mathfrak s_{\mathfrak d,k,\beta,\vec P}$
位于原仓西映射的
$\epsilon^{\mathcal K-(k+\|\beta\|)}\delta$
邻域内；扰动模空间
$\mathcal M_{k+1}(L;k,\beta;\vec P)
^{\mathfrak s_{\mathfrak d,k,\beta,\vec P}}$
位于原仓西映射零集的同一邻域内。

此处的
$\epsilon^{\mathcal K-(k+\|\beta\|)}\delta$
与 $P_i$ 的代无关；
而 (7.2.56.1) 中的
$\epsilon^{C(\mathcal K)-I}\delta$
还依赖于 $P_i$ 的代。
按此方式归纳也能成立，
其核验留给读者。
于是可对 $g=0,1,\ldots$
将全部 $\mathcal X_g(L)$ 取并，
得到一个可数生成的子复形，
$A_{n,K}$ 运算在其上处处有定义。
采用这一路径时，
$\mathfrak m_{k,\beta}$
与 $\mathfrak m^{\mathrm{geo}}_{k,\beta}$ 相同。
\end{remark}

本小节最后指出，
命题~\ref{prop:7-2-78}
立即推出定理 3.4.8 的如下变体。
\begin{corollary}{7.2.89}\label{cor:7-2-89}
\textup{（无滤过版本）}
当 $\beta=\beta_0$ 时，
$\mathfrak m_{k,\beta_0}$
在 $C_{(g_1)}(L;\mathbb Z)$ 上定义一个
$A_K$ 结构。
\end{corollary}
\begin{proof}
可同样证明命题~\ref{prop:7-2-35}
的一个只扰动
$\mathcal M^{\mathrm{main},\mathfrak d}_{k+1}(\beta_0)$
的版本，得到 $\mathcal X_g(L)$。
置
$\mathcal X_L=\bigcup_{g=0}^{K}\mathcal X_g(L)$。
容易检查 (3.4.6.1)--(3.4.6.4)。
这一步可用 $\mathbb Z$ 系数，
不必用 $\mathbb Q$ 系数：
因为只用到
$\mathcal M^{\mathrm{main},\mathfrak d}_{k+1}(\beta_0)$
的分支，所以可取单值截面，
不必取多值截面。
于是命题~\ref{prop:7-2-78} 的证明
也给出 $\mathbb Z$ 上的 $A_K$ 结构。
\end{proof}

\begin{remark}{7.2.90}\label{rem:7-2-90}
上述推论~\ref{cor:7-2-89} 的证明
与第~\ref{subsec:7-2-1} 小节的论证
看起来相当不同；这里说明二者的联系。
考虑由常值映射组成的模空间
$\mathcal M^{\mathrm{main},\mathfrak d}_{k+1}(\beta_0)$。
它与 $L\times Gr_{k+1}$ 微分同胚，
并且是横截的，即具有正确维数。
然而，其评价映射
$\mathcal M^{\mathrm{main},\mathfrak d}_{k+1}
(\beta_0)\to L^{k+1}$
并非浸没；
该映射在 $L$ 因子上是对角嵌入
$L\hookrightarrow L^{k+1}$，
在第二因子 $Gr_{k+1}$ 上平凡。
因而把障碍丛 $E$ 取为平凡丛所得的
显然仓西结构，
并不是命题 7.1.1、7.1.2 构造的结构。

后两个命题构造的仓西邻域，
可以取为 $L$ 在 $L^{k+1}$ 中的邻域。
扰动仓西映射 $\mathfrak s$ 后，
其零集 $\mathfrak s^{-1}(0)$
便把对角线 $L$ 扰动为
$L^{k+1}$ 中靠近 $L$ 的另一流形或链。
例如 $k=2$ 时，
$L\subset L^3$ 的扰动可写为
\[
\{(x,\varphi_1^{-1}(x),
\varphi_2^{-1}(x))\mid x\in L\}.
\]
再将 $L$ 与 $P_1\times P_2$ 作纤维积，
就得到第~\ref{subsec:7-2-1} 小节描述的
$\mathfrak m_2(P_1,P_2)$。
\end{remark}
